% !TEX TS-program = pdflatexmk
\RequirePackage{iftex}
\ifPDFTeX
  \pdfoutput=1
  \pdfminorversion=7
  \input glyphtounicode
  \pdfgentounicode=1
\fi
\ifdefined\pdfinterwordspaceon\pdfinterwordspaceon\fi
\documentclass[justified,notoc,twoside,nobib,nofonts]{tufte-handout}
\usepackage[T1]{fontenc}
\usepackage[utf8]{inputenc}
\usepackage[english]{babel}
\ifPDFTeX
  \usepackage[tracking]{microtype}
\else
  \usepackage{microtype}
\fi
\usepackage{amsthm,thmtools,thm-restate}
\usepackage{amsmath,amssymb,mathtools}
% Text and mathematics.
\usepackage[nomath,oldstylenums]{kpfonts}
\edef\KPbodyfamily{\rmdefault}
\renewcommand{\rmdefault}{zpltlf}
\usepackage{newpxmath}
\let\rmdefault\KPbodyfamily
\renewcommand{\familydefault}{\rmdefault}

% BaskervilleF headings use native Type 1 fonts and work with pdfLaTeX.
\newcommand{\BaskervilleHeading}{%
  \fontencoding{T1}\fontfamily{BaskervilleF-LF}%
  \fontseries{m}\fontshape{n}\selectfont}
\titleformat{\section}[hang]
  {\BaskervilleHeading\fontsize{15}{18}\selectfont}
  {\thesection}{1em}{}
\titleformat{\subsection}[hang]
  {\BaskervilleHeading\fontsize{13}{16}\selectfont}
  {\thesubsection}{1em}{}
\titleformat{\paragraph}[runin]
  {\BaskervilleHeading}{\theparagraph}{1em}{}
\usepackage{etoolbox}
\makeatletter
\patchcmd{\maketitle}{\LARGE\itshape\@title}
  {\BaskervilleHeading\fontsize{22}{26}\selectfont\raggedright\@title}
  {}{\PackageError{unique-factorisation-snarks}{Title font patch failed}{}}
% The author and email are printed immediately below the title block.
\patchcmd{\maketitle}{\@@author\@@date}{\@@date}
  {}{\PackageError{unique-factorisation-snarks}{Author block patch failed}{}}
\makeatother
\usepackage{booktabs,array,enumitem,needspace}
\usepackage[numbers,square]{natbib}
\usepackage{xcolor}
\usepackage{xurl,hyperref,cleveref}
\urlstyle{same}
\geometry{a4paper,top=23mm,bottom=46mm,left=28mm,right=56mm,marginparwidth=3cm}
\setlength{\emergencystretch}{2em}
\clubpenalty=10000
\widowpenalty=10000
\displaywidowpenalty=10000
\setcounter{secnumdepth}{2}
\AtBeginDocument{\fontsize{12}{15}\selectfont}
\fancyhead[LE]{\thepage\quad\textsmallcaps{nikolay ulyanov}}
\fancyhead[RO]{\textsmallcaps{factorisations of snarks}\quad\thepage}
\declaretheorem[name=Theorem,numberwithin=section]{theorem}
\declaretheorem[name=Proposition,sibling=theorem]{proposition}
\declaretheorem[name=Lemma,sibling=theorem]{lemma}
\declaretheorem[name=Corollary,sibling=theorem]{corollary}
\declaretheorem[name=Definition,sibling=theorem,style=definition]{definition}
\setlist[enumerate]{label=\textup{(\roman*)},leftmargin=*,itemsep=2pt,topsep=4pt}
\crefname{theorem}{theorem}{theorems}
\crefname{lemma}{lemma}{lemmas}
\crefname{definition}{definition}{definitions}
\setcitestyle{numbers,square}

% Full references appear in the margin at their first citation.
% Keep the canonical margin-selection block below when cloning this paper.
\makeatletter
\newdimen\@gp@fnmarkwidth
\newcommand{\bothcite}[2][0cm]{%
  \citep{#2}\sidenote[][#1]{\raggedright\csname paper@#2\endcsname}%
  \kern\@gp@fnmarkwidth\@gp@cite@space}
\def\@gp@cite@space{\@ifnextchar.{\@gp@cite@halfspace}{\@gp@cite@comma}}
\def\@gp@cite@comma{\@ifnextchar,{\@gp@cite@halfspace}{\@gp@cite@semicolon}}
\def\@gp@cite@semicolon{\@ifnextchar;{\@gp@cite@halfspace}{\@gp@cite@colon}}
\def\@gp@cite@colon{\@ifnextchar:{\@gp@cite@halfspace}{\@gp@cite@exclamation}}
\def\@gp@cite@exclamation{\@ifnextchar!{\@gp@cite@halfspace}{\@gp@cite@question}}
\def\@gp@cite@question{\@ifnextchar?{\@gp@cite@halfspace}{\@gp@cite@openparen}}
\def\@gp@cite@openparen{\@ifnextchar({\@gp@cite@halfspace}{\@gp@cite@closeparen}}
\def\@gp@cite@closeparen{\@ifnextchar){\@gp@cite@halfspace}{\@gp@cite@bracket}}
\def\@gp@cite@bracket{\@ifnextchar]{\@gp@cite@halfspace}{\@gp@cite@normalspace}}
\def\@gp@cite@halfspace{\nobreak\hskip.5\fontdimen2\font\relax}
\def\@gp@cite@normalspace{\nobreak\space}
% GP margin selection v1 (pdfLaTeX). Keep this block in the standalone source.
\makeatletter
\@ifundefined{@gp@fnmarkwidth}{\newdimen\@gp@fnmarkwidth}{}
\ifPDFTeX
\def\@makefnmark{%
  \hbox{%
    \setbox\z@=\hbox{\@textsuperscript{\normalfont\footnotesize\@thefnmark}}%
    \global\@gp@fnmarkwidth=\wd\z@%
    \setbox\tw@=\copy\z@%
    \immediate\pdfxform\tw@%
    \pdfannot width\wd\z@ height\ht\z@ depth\dp\z@{%
      /Subtype /Watermark /F 708 /Contents ()%
      /AP << /N \the\pdflastxform\space 0 R >>%
    }%
  }%
}
\fi
% Selection-safe margin text: annotation appearances, with separate URI links.
\ifPDFTeX
\newbox\@gp@marginbox
\newdimen\@gp@marginwidth
\newdimen\@gp@marginheight
\newdimen\@gp@marginfirstheight
\newdimen\@gp@marginpadding
\newdimen\@gp@marginskip
\newdimen\@gp@margindepth
\newcount\@gp@marginserial
\long\def\@gp@marginappearance#1{%
  \begingroup
  \global\let\@gp@marginuri\@empty
  \setbox\@gp@marginbox=\vtop{%
    \hsize=\marginparwidth
    \def\href##1##2{%
      \ifx\@gp@marginuri\@empty
        \xdef\@gp@marginuri{\detokenize{##1}}%
      \else
        \PackageError{gp-margin-selection}{More than one URI in a margin note}%
          {Use a separate note for each reference.}%
      \fi
      \textcolor{\@urlcolor}{##2}%
    }%
    #1\par
  }%
  \@gp@marginfirstheight=\ht\@gp@marginbox
  \setbox\@gp@marginbox=\vbox{\unvbox\@gp@marginbox}%
  \@gp@marginpadding=\baselineskip
  \advance\@gp@marginpadding-\prevdepth
  \advance\@gp@marginpadding-\@gp@marginfirstheight
  \ifdim\@gp@marginpadding<\lineskiplimit
    \@gp@marginpadding=\lineskip
  \fi
  \@gp@marginskip=\baselineskip
  \advance\@gp@marginskip-\prevdepth
  \advance\@gp@marginskip-\ht\@gp@marginbox
  \ifdim\@gp@marginskip<\lineskiplimit
    \@gp@marginskip=\lineskip
  \fi
  \advance\@gp@marginpadding-\@gp@marginskip
  \vskip\@gp@marginpadding
  \@gp@marginwidth=\wd\@gp@marginbox
  \@gp@marginheight=\ht\@gp@marginbox
  \@gp@margindepth=\dp\@gp@marginbox
  \immediate\pdfxform\@gp@marginbox
  \edef\@gp@marginform{\the\pdflastxform}%
  \global\advance\@gp@marginserial\@ne
  \hbox{%
    \pdfannot width\@gp@marginwidth height\@gp@marginheight depth\@gp@margindepth{%
      /Subtype /Watermark /F 708 /Contents ()%
      /NM (GPmargin-\the\@gp@marginserial)%
      /AP << /N \@gp@marginform\space 0 R >>%
    }%
    \ifx\@gp@marginuri\@empty\else
      \pdfannot width\@gp@marginwidth height\@gp@marginheight depth\@gp@margindepth{%
        /Subtype /Link /F 4 /Border [0 0 0]%
        /A << /S /URI /URI (\pdfescapestring{\@gp@marginuri}) >>%
      }%
    \fi
    \vrule width0pt height\@gp@marginheight depth\@gp@margindepth
    \kern\@gp@marginwidth
  }%
  \endgroup
}
\else
\long\def\@gp@marginappearance#1{#1}
\fi
\renewcommand\@footnotetext[2][0pt]{%
  \marginpar{%
    \hbox{}\vspace*{#1}%
    \def\baselinestretch{\setspace@singlespace}%
    \reset@font\footnotesize%
    \@tufte@margin@par%
    \vspace*{-1\baselineskip}%
    \@gp@marginappearance{%
      \noindent
      \protected@edef\@currentlabel{%
        \csname p@footnote\endcsname\@thefnmark%
      }%
      \color@begingroup
        \@makefntext{\ignorespaces#2}%
      \color@endgroup
    }%
  }%
}
\renewcommand\marginnote[2][0pt]{%
  \ifthenelse{\boolean{@tufte@loadnatbib}}{%
    \let\cite\@tufte@infootnote@cite%
  }{}%
  \gdef\@tufte@citations{}%
  \marginpar{%
    \hbox{}\vspace*{#1}%
    \@tufte@marginnote@font\@tufte@marginnote@justification%
    \@tufte@margin@par\vspace*{-1\baselineskip}%
    \@gp@marginappearance{\noindent #2}%
  }%
  \@tufte@print@citations%
  \ifthenelse{\boolean{@tufte@loadnatbib}}{%
    \let\cite\@tufte@normal@cite%
  }{}%
}
\makeatother

\expandafter\def\csname paper@CS10\endcsname{%
  Miroslav Chladn\'y and Martin \v{S}koviera.
  \emph{Factorisation of snarks}.
  Electronic Journal of Combinatorics 17(1) (2010), R32.
  \href{https://doi.org/10.37236/304}{\nolinkurl{doi:10.37236/304}}.}
\expandafter\def\csname paper@GZ18\endcsname{%
  Jan Goedgebeur and Carol T. Zamfirescu.
  \emph{On hypohamiltonian snarks and a theorem of Fiorini}.
  Ars Mathematica Contemporanea 14(2) (2018), 227--249.
  \href{https://doi.org/10.26493/1855-3974.1176.9b7}{\nolinkurl{doi:10.26493/1855-3974.1176.9b7}}.}
\expandafter\def\csname paper@MS19\endcsname{%
  Edita M\'a\v{c}ajov\'a and Martin \v{S}koviera.
  \emph{Permutation snarks of order $2\pmod 8$}.
  Acta Mathematica Universitatis Comenianae 88(3) (2019), 929--934.
  \href{https://www.iam.fmph.uniba.sk/amuc/ojs/index.php/amuc/article/view/1313}{journal version}.}
\expandafter\def\csname paper@MS25\endcsname{%
  Edita M\'a\v{c}ajov\'a and Martin \v{S}koviera.
  \emph{Are there any permutation snarks on $6\pmod 8$ vertices?}
  Procedia Computer Science 273 (2025), 163--168.
  \href{https://doi.org/10.1016/j.procs.2025.10.294}{\nolinkurl{doi:10.1016/j.procs.2025.10.294}}.}
\expandafter\def\csname paper@BGHM13\endcsname{%
  Gunnar Brinkmann, Jan Goedgebeur, Jonas H\"agglund and Klas Markstr\"om.
  \emph{Generation and properties of snarks}.
  Journal of Combinatorial Theory, Series B 103(4) (2013), 468--488.
  \href{https://doi.org/10.1016/j.jctb.2013.05.001}{\nolinkurl{doi:10.1016/j.jctb.2013.05.001}}.}
\expandafter\def\csname paper@NS95\endcsname{%
  Roman Nedela and Martin \v{S}koviera.
  \emph{Atoms of cyclic connectivity in cubic graphs}.
  Mathematica Slovaca 45(5) (1995), 481--499.
  \href{https://dml.cz/dmlcz/132921}{DML-CZ:132921}.}
\newcommand{\paperreference}[1]{\csname paper@#1\endcsname}
\makeatother

\title[Graph Puzzles IV.1-preview: Factorisations of snarks]{Graph Puzzles IV.1-preview: Factorisations of hypohamiltonian and permutation snarks}
\author{Nikolay Ulyanov + AI-Slop}
\date{6 October 2026}
\hypersetup{
  pdftitle={Graph Puzzles IV.1-preview: Factorisations of hypohamiltonian and permutation snarks},
  pdfauthor={Nikolay Ulyanov + AI-Slop},
  colorlinks=true,
  linkcolor=blue!50!black,
  citecolor=blue!50!black,
  urlcolor=blue!50!black
}

\begin{document}
\selectlanguage{english}
\maketitle
Nikolay Ulyanov\textsuperscript{*} + AI-Slop%
\marginnote{\textsuperscript{*}{\scriptsize\href{mailto:ulyanick@gmail.com}{ulyanick@gmail.com}}}\par

\begin{abstract}
\noindent\newthought{Abstract: }
Every hypohamiltonian snark and every permutation snark has a unique
multiset of cyclically $5$-edge-connected factors of the same class under
successive decompositions along cycle-separating $4$-edge-cuts.
We prove that hypohamiltonicity passes from a dot product to both of its
snark factors. Together with the factorisation theorem for bicritical
snarks, this answers a question of Goedgebeur and Zamfirescu.
For permutation snarks, we extend the atom argument of Chladn\'y and
\v{S}koviera using colourable cut multipoles and closure under their
canonical completions. The two conclusions also hold within the class
of hypohamiltonian permutation snarks.
\end{abstract}
\noindent\rule{\linewidth}{0.4pt}

\Needspace{12\baselineskip}
\section{Introduction}

The dot product combines two snarks across a $4$-edge-cut. Reversing this
operation repeatedly expresses a snark in terms of factors with higher
cyclic connectivity. Chladn\'y and \v{S}koviera~\bothcite{CS10}
proved that a bicritical snark has a unique multiset of cyclically
$5$-edge-connected bicritical factors. Their theorem identifies the
isomorphism types of the factors and their multiplicities, independently
of the cuts chosen during the decomposition.

A graph is \emph{hypohamiltonian} if it is non-Hamiltonian and every
single-vertex deletion is Hamiltonian. Hypohamiltonian snarks are
bicritical. Goedgebeur and Zamfirescu~\bothcite{GZ18}
asked whether every hypohamiltonian snark of cyclic connectivity four
is a dot product of two hypohamiltonian snarks
\citep[Problem~4.3]{GZ18}. We prove that both snark factors of any
hypohamiltonian dot product are hypohamiltonian. The unique
factorisation theorem for bicritical snarks then gives the following
result.

\begin{restatable}{theorem}{hypofactorisation}\label{thm:hypo}
Every hypohamiltonian snark has a unique multiset, up to isomorphism and
with multiplicities, of cyclically $5$-edge-connected hypohamiltonian
snark factors. Every intermediate factor in a decomposition along
cycle-separating $4$-edge-cuts is hypohamiltonian.
\end{restatable}

A \emph{permutation $2$-factor} consists of two induced cycles. A snark
admitting such a $2$-factor is a \emph{permutation snark}. Its two cycles
have equal odd length, and their complement is a perfect matching
between them. M\'a\v{c}ajov\'a and \v{S}koviera~\bothcite{MS19}
proved that every cycle-separating $4$-edge-cut of a permutation snark
determines two smaller permutation snarks
\citep[Section~2 and Theorem~2.1]{MS19}.

Permutation snarks are conjectured to be bicritical; see
M\'a\v{c}ajov\'a and \v{S}koviera~\bothcite{MS25}, Section~1.
Together with the factorisation theorem for bicritical snarks,
this conjecture makes it natural to ask whether permutation snarks
also have unique factorisations. We use the closure property above
to prove the following result.

\begin{restatable}{theorem}{permfactorisation}\label{thm:perm}
Every permutation snark has a unique multiset, up to isomorphism and
with multiplicities, of cyclically $5$-edge-connected permutation
snark factors. Every intermediate factor in a decomposition along
cycle-separating $4$-edge-cuts is a permutation snark.
\end{restatable}

Permutation snarks include non-hypohamiltonian graphs. There is a
$26$-vertex example obtained as a successive dot product of three
Petersen graphs, and there is also a $50$-vertex example with cyclic
connectivity five.
Brinkmann, Goedgebeur, H\"agglund and Markstr\"om~\bothcite{BGHM13}
enumerated $64$ permutation snarks on $26$ vertices
\citep[Table~3]{BGHM13}.
Thus the hypohamiltonian permutation snarks form a proper subclass
of the permutation snarks, also within the cyclically
$5$-edge-connected class. The two factorisation theorems apply to
the same completed factors at each cut, giving the following result
for their intersection.

\begin{restatable}{corollary}{intersectionfactorisation}\label{cor:intersection}
Every hypohamiltonian permutation snark has a unique multiset, up to
isomorphism and with multiplicities, of cyclically $5$-edge-connected
hypohamiltonian permutation snark factors. Every intermediate factor
has both properties.
\end{restatable}

\Needspace{12\baselineskip}
\section{Dot products and cut multipoles}\label{sec:preliminaries}

All graphs are finite. We use \emph{snark} for a simple cubic graph of
girth at least five, cyclically $4$-edge-connected, and with chromatic
index four. A \emph{cycle} is a connected $2$-regular subgraph.
For $X\subseteq V(G)$, write $\partial_G X$ for the edges with exactly
one end in $X$. The cut is \emph{cycle-separating} when both $G[X]$ and
$G[V(G)\setminus X]$ contain a cycle. We write $\lambda_c(G)$ for
cyclic edge-connectivity, $g(G)$ for girth, and $|G|$ for $|V(G)|$.

\subsection{The dot product}

Let $A$ and $B$ be disjoint snarks. Choose independent edges $ab,cd$ of
$A$ and an edge $xy$ of $B$, with
\[
 N_B(x)=\{a',b',y\},\qquad N_B(y)=\{c',d',x\}.
\]
The four vertices $a',b',c',d'$ are distinct. Put
\[
 L=A-\{ab,cd\},\qquad R=B-\{x,y\},
\]
where edges are deleted from $A$ and vertices from $B$. The dot product
with this attachment is
\begin{equation}\label{eq:dot-product}
 G=A\mathbin{\cdot}B
   =L\cup R+\{aa',bb',cc',dd'\}.
\end{equation}
The four joining edges form the \emph{bond} $F$. Its two couples are
$\{aa',bb'\}$ and $\{cc',dd'\}$. This is the convention of
\citep[Section~2]{CS10} and \citep[Section~2]{GZ18}.
The order satisfies
\begin{equation}\label{eq:order}
 |G|-2=(|A|-2)+(|B|-2).
\end{equation}

A cycle-separating $4$-edge-cut in a snark is a matching. Indeed, if a
vertex is incident with at least two cut edges, moving it to the other
shore reduces the size of the cut and preserves a cycle in each shore.
The resulting cut has size at most three, contrary to cyclic
$4$-edge-connectivity. Thus each shore has four distinct boundary
vertices.

\subsection{Canonical completions}

Disconnecting the four edges of a cut leaves a \emph{$4$-pole} on each
shore: a cubic multipole with four semiedges. A proper $3$-edge-colouring
of a multipole assigns colours to its ordinary edges and semiedges,
with three distinct colours incident at each vertex. We use the
multipole terminology of \citep[Section~3]{CS10}.

Up to a permutation of colours, the Parity Lemma permits four types
of boundary colouring:
\[
 1111,\qquad1122,\qquad1212,\qquad1221.
\]
Every colourable $4$-pole admits at least two of these types. If two
colourable $4$-poles join to a snark, their sets of types are disjoint,
so each admits exactly two. Such a pole is \emph{colour-open}.
Its four semiedges have a unique partition into two couples. In the
\emph{isochromatic} case, the semiedges of each couple receive the same
colour in every colouring. In the \emph{heterochromatic} case, the
semiedges of each couple receive distinct colours in every colouring.
The two poles of the cut have opposite types, and their couples agree
across the cut \citep[Propositions~3.1--3.5]{CS10}.

The \emph{canonical snark completion} of a heterochromatic pole joins
the semiedges within each couple. For an isochromatic pole, it attaches
each couple to a new vertex and joins the two new vertices. These
completions are unique up to isomorphism. They recover the two factors
in \eqref{eq:dot-product}, with the heterochromatic completion playing
the role of $A$ and the isochromatic completion the role of $B$.

A \emph{composition chain} of $G$ is a multiset obtained from $\{G\}$
by successively replacing a member with its two canonical factors
at a cycle-separating $4$-edge-cut. A chain is \emph{terminal} when every
member is cyclically $5$-edge-connected. Two chains are
\emph{equivalent} if their members can be matched, including
multiplicities, by graph isomorphisms. We use the terminology of
\citep[Section~7]{CS10}.

\Needspace{12\baselineskip}
\section{Hypohamiltonian factors}\label{sec:hypohamiltonian}

We first record the boundary restrictions imposed by the
non-Hamiltonicity of the two factors.

\Needspace{9\baselineskip}
\begin{lemma}\label{lem:boundary}
In \eqref{eq:dot-product}, let $A$ and $B$ be non-Hamiltonian graphs.
A spanning path in $L$ with terminal ends joins different pairs among
$\{a,b\}$ and $\{c,d\}$. A spanning path in $R$ with terminal ends
joins $a'$ to $b'$ or $c'$ to $d'$.
Two disjoint paths spanning $L$, with ends $a,b,c,d$, pair $a$ with $b$
and $c$ with $d$. Two disjoint paths spanning $R$, with ends
$a',b',c',d'$, pair $a'$ with $b'$ and $c'$ with $d'$.
\end{lemma}
\begin{proof}
A spanning $a$--$b$ path in $L$ closes to a Hamiltonian cycle of $A$
when $ab$ is restored; the same holds for $c,d$. A spanning path in
$R$ whose ends belong to different pairs closes through $x,y$ to a
Hamiltonian cycle of $B$.

For a spanning pair of paths in $L$, either cross pairing of the four
ends becomes one Hamiltonian cycle after restoring $ab$ and $cd$.
For a spanning pair in $R$, either cross pairing becomes one
Hamiltonian cycle after adjoining $a'xb'$ and $c'yd'$.
The remaining pairing is the one stated in each case.
\end{proof}

\Needspace{9\baselineskip}
\begin{lemma}\label{lem:descent}
If a dot product $G=A\mathbin{\cdot}B$ of snarks is hypohamiltonian,
then both $A$ and $B$ are hypohamiltonian.
\end{lemma}
\begin{proof}
Use the notation of \eqref{eq:dot-product}. A Hamiltonian cubic graph
is $3$-edge-colourable, so both factors are non-Hamiltonian.
A Hamiltonian cycle in any vertex-deleted subgraph of $G$ meets $F$
in two or four edges. With two crossings, its restriction to the
undeleted shore is a spanning path; with four, that restriction
consists of two disjoint spanning paths.

Fix $v\in V(A)$ and take a Hamiltonian cycle $C$ in $G-v$.
If $C$ uses two bond edges, \cref{lem:boundary} shows that its spanning
path in $R$ has ends $a',b'$ or $c',d'$. Replace the whole segment
through $R$, including its bond edges, with $ab$ or $cd$, respectively.
If $C$ uses all four bond edges, its two spanning paths in $R$ pair
$a'$ with $b'$ and $c'$ with $d'$. Replace these two segments with
$ab$ and $cd$. Each replacement preserves the cyclic order of the
remaining vertices of $C$. The resulting cycle spans $A-v$.
Thus $A-v$ is Hamiltonian for every $v\in V(A)$.

Now fix $w\in V(B)\setminus\{x,y\}$ and take a Hamiltonian cycle in
$G-w$. With two bond edges, the spanning path in $L$ joins different
terminal pairs. For example, a path from $a$ to $c$ is replaced,
together with its bond edges, by $a'xyc'$. The other choices give the
corresponding path through $x,y$.
With four bond edges, the two spanning paths in $L$ pair $a$ with $b$
and $c$ with $d$. Replace them, together with their bond edges, by
$a'xb'$ and $c'yd'$. The resulting cycle spans $B-w$ in either case.

It remains to recover Hamiltonian cycles after deleting $x$ or $y$.
Take a Hamiltonian cycle of $G-a$. The missing bond edge $aa'$ forces
exactly two crossings. Its spanning path in $R$ has ends $c',d'$ by
\cref{lem:boundary}. Adjoining $c'yd'$ gives a Hamiltonian cycle of
$B-x$. A Hamiltonian cycle of $G-c$ similarly has a spanning
$a'$--$b'$ path in $R$; adjoining $a'xb'$ gives a Hamiltonian cycle
of $B-y$. Hence both factors are hypohamiltonian.
\end{proof}

A snark is \emph{bicritical} if deleting any two distinct vertices
yields a $3$-edge-colourable graph. Hypohamiltonian snarks have this
property \citep[Section~1]{GZ18}. For distinct $u,v$, deleting $v$
from a Hamiltonian cycle of $G-u$ leaves a spanning path in
$G-\{u,v\}$. Colour the path alternately with colours $1,2$.
The remaining edges form a matching and receive colour $3$.

\Needspace{11\baselineskip}
\hypofactorisation*
\begin{proof}
Let $G$ be hypohamiltonian and let $S$ be any cycle-separating
$4$-edge-cut. By bicriticality and \citep[Theorem~7.1]{CS10},
$S$ determines a dot product $G=A\mathbin{\cdot}B$ of smaller
bicritical snarks. By \cref{lem:descent}, both factors are
hypohamiltonian. Repeating the argument keeps every factor
hypohamiltonian.

Each decomposition replaces a graph with two smaller graphs, so
every chain has a terminal refinement. Its members are cyclically
$5$-edge-connected hypohamiltonian snarks. The bicritical
unique-factorisation theorem \citep[Theorem~C, or Theorem~10.3]{CS10}
identifies any two such terminal chains up to isomorphism and order,
including multiplicities.
\end{proof}

In particular, every hypohamiltonian snark of cyclic connectivity four
is a dot product of two hypohamiltonian snarks. This gives an
affirmative answer to \citep[Problem~4.3]{GZ18}.

\Needspace{12\baselineskip}
\section{A factorisation criterion}\label{sec:criterion}

The permutation result uses the atom argument for bicritical snarks
\citep[Sections~8--10]{CS10}. We formulate the argument with the
closure hypotheses that enter its proof.

Let $\mathcal C$ be a class of snarks closed under isomorphism.
Assume that every cycle-separating $4$-edge-cut in a member of
$\mathcal C$ has two colourable shore poles, and that both canonical
completions belong to $\mathcal C$.
For a cyclic fragment $X$, write $\widetilde X$ for the completion of
its shore pole and $G/\widetilde X$ for the completion of the
complementary pole.

Suppose $\lambda_c(G)=4$. A \emph{cyclic fragment} is an induced
shore of a cycle-separating $4$-edge-cut, and an \emph{atom} is a
cyclic fragment minimal under inclusion. Since $4<g(G)$, the atom
theorem of Nedela and \v{S}koviera~\bothcite{NS95}
places every atom wholly on one side of every cycle-separating
$4$-edge-cut; see also \citep[Theorem~7.2]{CS10}.
The cut $\partial_G X$ of an atom is \emph{atomic}.
A cycle-separating cut with a shore $X\cup\{u,v\}$, where $u,v$ are
adjacent outside vertices each joined once to $X$, is
\emph{quasiatomic}. These are the cuts associated with $X$.

\Needspace{9\baselineskip}
\begin{lemma}\label{lem:atom-switching}
Let $G\in\mathcal C$ have cyclic connectivity four and let $X$ be an
atom. A cut associated with $X$ gives factors isomorphic to
$\widetilde X$ and $G/\widetilde X$.
For any other cycle-separating $4$-edge-cut $S$, one $S$-factor $P$
contains $X$ as an atom, and a corresponding cut $p_X(S)$ is
cycle-separating in $G/\widetilde X$.
The two refinements obtained by using $S$ and the atomic cut in
either order give equivalent three-member chains.
\end{lemma}
\begin{proof}
We follow the cut analysis of \citep[Proposition~8.2, Lemmas~9.1--9.2,
Propositions~9.3--9.4, and Lemma~10.1]{CS10}. The hypotheses above
provide colour-open poles and their completions at every step;
each completed graph has girth at least five and cyclic connectivity
at least four.

First consider a quasiatomic shore $X\cup\{u,v\}$.
Its pole is heterochromatic: an isochromatic completion would
place the adjacent vertices $u,v$ beside the two new cap vertices
and create a triangle or a quadrilateral. Joining its boundary
couples instead completes the pole of $X$ by an adjacent pair.
Thus the completed shore is isomorphic to $\widetilde X$.
In the complementary completion, the same isomorphism exchanges the
original pair $u,v$ with the added pair and fixes the other vertices.
The associated cuts therefore give equivalent factor pairs.

For the remaining cuts, put $T=\partial_G X$.
The following two restrictions control the completions:
\begin{equation}\label{eq:overlap}
 |S\cap T|\ne3,\qquad
 |S\cap T|=2\ \Longrightarrow\
 S\cap T\text{ is a couple of neither cut}.
\end{equation}
For three common edges, the two remaining edges of $S\cup T$
separate a cyclic $2$-pole from $X$, contrary to cyclic
$4$-edge-connectivity. Here the $2$-pole has vertices, so the cubic
degree count forces it to contain a cycle.
For two common edges forming an $S$-couple, complete the $S$-shore
containing $X$. A heterochromatic completion produces a
cycle-separating $2$-edge-cut; an isochromatic completion produces an
independent $3$-edge-cut. Both contradict cyclic
$4$-edge-connectivity of the completed factor.
If the two common edges form a $T$-couple, restricting colourings of
the larger shore pole to the pole of $X$ transfers the corresponding
equality or inequality of colours to them. They consequently
form an $S$-couple as well. This proves \eqref{eq:overlap}.

The atom theorem puts $X$ inside one $S$-shore. Its completion $P$
keeps $X$ induced: an added edge within $X$ would join the two common
edges as an $S$-couple, contrary to \eqref{eq:overlap}.
If $P-X$ is acyclic, the cubic $4$-pole degree count leaves two
adjacent vertices; reversing the completion identifies $S$ as an
atomic or quasiatomic cut. For the cuts under consideration,
$P-X$ therefore contains a cycle. Thus $X$ remains a cyclic fragment
of $P$. Its minimality follows because every cyclic fragment
properly contained in $X$ in $P$ would also be a cyclic fragment of
$G$. Hence $X$ is an atom of $P$.

To describe the corresponding cut, replace the pole of $X$ with
its colour-equivalent acyclic $4$-pole. For a heterochromatic pole,
this acyclic pole consists of two adjacent vertices with two
semiedges at each; for an isochromatic pole, it consists of two
isolated edges. The resulting graph is $G/\widetilde X$.
The map $p_X$ keeps edges outside $X$ and maps the four boundary
edges to their replacement edges. In the second case, the two
edges of each couple have the same image.
By \eqref{eq:overlap}, $p_X(S)$ has four distinct edges.
The shore of $G-S$ containing $X$ also has vertices outside $X$,
since $S\ne\partial_G X$. After deleting $p_X(S)$, these vertices
remain separated from the other shore. Thus $p_X(S)$ is an edge-cut.

The images are also independent. In the adjacent-vertex replacement,
an adjacency could arise only from a boundary couple contained in
$S$, which \eqref{eq:overlap} excludes. In the isolated-edge
replacement, an inherited edge might become adjacent to a new
joined edge. Move their common vertex to the other shore of the
projected cut. This gives a $3$-edge-cut, which isolates a vertex in
the completed snark. The other two projected edges therefore meet
there. Their preimages in $S$ are independent, so one is a second
boundary edge of $X$. The resulting configuration consists of $X$
and the adjacent pair of these two outside vertices; its boundary
is $S$. This is precisely the quasiatomic configuration.
Thus $p_X(S)$ is independent. Both shores of an independent cut in
a cubic graph have minimum degree at least two, and hence contain
cycles. This proves the asserted projection property.

Finally, let $P,Q$ be the $S$-factors, with $X\subseteq P$.
Decomposing $P$ at $X$ gives the chain
\[
 \{\widetilde X,\ P/\widetilde X,\ Q\}.
\]
If the atomic cut is used first, decomposing $G/\widetilde X$ at
$p_X(S)$ leaves the pole defining $Q$ intact. The other pole is
obtained from the pole defining $P$ by the colour-equivalent
replacement of $X$. Colour-equivalent substitution preserves its
boundary colourings, so its canonical completion is
$P/\widetilde X$. The two chains are equivalent.
\end{proof}

\Needspace{9\baselineskip}
\begin{lemma}\label{lem:splitting}
For $G\in\mathcal C$ of cyclic connectivity four and an atom $X$,
every terminal chain of $G$ is equivalent to the union of a terminal
chain of $\widetilde X$ and a terminal chain of $G/\widetilde X$.
\end{lemma}
\begin{proof}
Follow $X$ through a terminal chain. By \cref{lem:atom-switching},
each cut outside the associated cuts preserves $X$ as an atom in
one factor. Eventually an associated cut is used, since a terminal
factor has cyclic connectivity at least five.
Replace that step with the atomic cut using the first assertion of
\cref{lem:atom-switching}.

Move the atomic step towards the start of the chain. It commutes
with a preceding step in a different factor. When the preceding
step created the factor containing $X$, the last assertion of
\cref{lem:atom-switching} interchanges these two steps and preserves
the resulting three factors up to isomorphism. Transport the later
refinements across these isomorphisms. Each interchange moves the
atomic step one position earlier, so finitely many interchanges
place it first. The subsequent steps refine its two factors,
yielding the two terminal chains in the statement.
This is the splitting argument of \citep[Lemma~10.2]{CS10}.
\end{proof}

\Needspace{11\baselineskip}
\begin{theorem}\label{thm:criterion}
Let $\mathcal C$ be a class of snarks closed under isomorphism.
Suppose that the shore poles of every cycle-separating $4$-edge-cut
in a member of $\mathcal C$ are colourable, and that their canonical
snark completions belong to $\mathcal C$.
Then every member of $\mathcal C$ has a unique multiset, up to
isomorphism and with multiplicities, of cyclically
$5$-edge-connected factors in $\mathcal C$.
\end{theorem}
\begin{proof}
A canonical decomposition has two factors in $\mathcal C$ satisfying
\eqref{eq:order}. Each has at least ten vertices, so both are smaller
than the parent. Thus every chain has a terminal refinement, and
the terminal members have cyclic connectivity at least five.

For uniqueness, use induction on $|G|$. A cyclically
$5$-edge-connected member has the single-member terminal chain
$\{G\}$. Otherwise choose an atom $X$ and let $\mathcal F$ and
$\mathcal F'$ be two terminal chains. By \cref{lem:splitting},
\[
 \begin{aligned}
 \mathcal F&\simeq\mathcal F_1\uplus\mathcal F_2,\\
 \mathcal F'&\simeq\mathcal F'_1\uplus\mathcal F'_2,
 \end{aligned}
\]
where the first subchains factor $\widetilde X$ and the second
factor $G/\widetilde X$. Both graphs are smaller members of
$\mathcal C$. The induction hypothesis gives
$\mathcal F_1\simeq\mathcal F'_1$ and
$\mathcal F_2\simeq\mathcal F'_2$, and hence
$\mathcal F\simeq\mathcal F'$.
\end{proof}

\Needspace{12\baselineskip}
\section{Permutation factors}\label{sec:permutation}

We verify the two closure hypotheses of \cref{thm:criterion} using
a permutation $2$-factor. The fixed-cut closure is the result of
\citep[Section~2]{MS19}; the following proof also describes its
colour-open completions.

\Needspace{9\baselineskip}
\begin{lemma}\label{lem:permutation-cut}
Let $G$ be a permutation snark with permutation $2$-factor
$C_1\cup C_2$, and let $S$ be a cycle-separating $4$-edge-cut.
The cut consists of two edges of each $C_i$. Each shore pole is
colourable, and its canonical completion is a permutation snark.
\end{lemma}
\begin{proof}
Write $M=E(G)\setminus E(C_1\cup C_2)$ and let the common length of
the cycles be $m$. Each cycle is cut by $S$.
Indeed, if a whole $C_i$ lies in one shore, every vertex in the
other shore sends its matching edge across $S$. That shore then
has at most four vertices, and girth at least five makes it
acyclic. This contradicts cycle separation.
Each cycle contributes a positive even number of cut edges.
Thus $S$ has two edges on each cycle, and every edge of $M$ stays
within a shore.

Each shore contains one path from each $C_i$. Its matching edges
pair their vertices, so the two paths have the same order $r$.
Colour the matching edges $3$ and alternate $1,2$ along each path,
including the two boundary semiedges at its ends.
The initial colour on each path can be chosen independently.
This gives a proper colouring of the shore pole.

The two path ends form the boundary couples. If $r$ is even, the
semiedges at the ends of each path receive the same colour.
The independent phase choices give types $1111$ and $1122$, with
the two paths defining the pairs. The pole is isochromatic and its
completion adds one vertex to close each path, together with a
matching edge between the new vertices.
If $r$ is odd, the end semiedges receive distinct colours, and the
phase choices give types $1212$ and $1221$. The pole is
heterochromatic and its completion joins the ends of each path.
In both cases the completed graph has two induced cycles and a
perfect matching between them. Its canonical completion has
chromatic index four. An uncolourable cycle permutation graph is
cyclically $4$-edge-connected and has girth at least five
\citep[Section~1]{MS19}; hence the completed graph is a permutation
snark.
\end{proof}

\Needspace{11\baselineskip}
\permfactorisation*
\begin{proof}
Apply \cref{thm:criterion} to the class of permutation snarks.
By \cref{lem:permutation-cut}, every cycle-separating $4$-edge-cut
has colourable shore poles, and both canonical factors are
permutation snarks. The theorem gives existence and uniqueness
of the terminal factor multiset. The same closure statement
applied at each step keeps all intermediate factors in the
permutation class.
\end{proof}

\Needspace{11\baselineskip}
\intersectionfactorisation*
\begin{proof}
At any cycle-separating $4$-edge-cut of a hypohamiltonian permutation
snark, \cref{lem:permutation-cut} gives two canonical permutation
snark factors. By \cref{lem:descent}, both are hypohamiltonian.
This applies at every later step. Existence and uniqueness follow
from \cref{thm:hypo,thm:perm}. In particular, the hypohamiltonian,
permutation and bicritical decompositions of this graph have the
same terminal multiset.
\end{proof}

\Needspace{12\baselineskip}
\section{Orders of composition factors}

The order identity \eqref{eq:order} gives a bound on the number of
terminal factors and a congruence for their orders.

\Needspace{12\baselineskip}
\begin{corollary}\label{cor:orders}
Let $H_1,\ldots,H_t$ be the terminal factors of a hypohamiltonian
or permutation snark $G$. Then
\[
 |G|-2=\sum_{i=1}^{t}(|H_i|-2),
 \qquad t\le\frac{|G|-2}{8}.
\]
For a permutation snark, each $|H_i|\equiv2\pmod4$ and
\[
 \frac{|G|-2}{4}\equiv
   \sum_{i=1}^{t}\frac{|H_i|-2}{4}\pmod2.
\]
\end{corollary}
\begin{proof}
Iterate \eqref{eq:order} along the decomposition. Every snark factor
has order at least ten, giving the bound on $t$.
For permutation factors, the two induced cycles have equal odd
length, so their orders are $2\pmod4$. Dividing the order identity
by four gives the congruence.
\end{proof}

Consequently the existence of a permutation snark of order
$6\pmod8$ is equivalent to the existence of a cyclically
$5$-edge-connected permutation snark of that order residue.
Indeed, an odd value of $(|G|-2)/4$ forces an odd contribution from
at least one terminal factor.

\Needspace{12\baselineskip}
\begingroup
\fontsize{10}{12}\selectfont
\setlength{\bibsep}{5pt}
\begin{thebibliography}{BGHM13}
\bibitem[CS10]{CS10}\paperreference{CS10}
\bibitem[GZ18]{GZ18}\paperreference{GZ18}
\bibitem[MS19]{MS19}\paperreference{MS19}
\bibitem[MS25]{MS25}\paperreference{MS25}
\bibitem[BGHM13]{BGHM13}\paperreference{BGHM13}
\bibitem[NS95]{NS95}\paperreference{NS95}
\end{thebibliography}
\endgroup
\end{document}
